% Faithful typesetting transcription of the formal CCC2021.14 publication, Sections2--4.
\section*{2 Preliminaries}
We consider two-party randomized protocols for computing functions of the form $f:\{0,1\}^n\times\{0,1\}^n\to\{0,1\}$, while using a defective source of randomness. Specifically, we consider sources of randomness that produce $\ell$-bit long strings having min-entropy at least $k$; that is, each outcome occurs with probability at most $2^{-k}$. Such sources are called $(\ell,k)$-sources.
We consider both the public-randomness model in which the parties have access to common (public) randomness, and the private-randomness model in which each party has its private source of randomness, which is independent of the randomness of the other party. In our context (of defective random sources) it is important to stress that the postulated sources of randomness are the only ones available to the parties.
The results hold not only for ``alternating protocols'' (in which the parties alternatively exchange single bits), but directly for any protocol in which the sender of the next bit is determined by the communication so far; that is, no need to lose a factor of two in translation (from such general protocols to ``alternating'' ones).
\subsection*{2.1 Specific Background About Communication Complexity}
We shall use the following basic result that refers to deterministic communication protocols.
\begin{authorclaim}{Claim2.1 (the ``corners lemma'' (cf., e.g., [7, Prop.1.13--1.14] or [10, Lem.1.3--1.4]))}
Let $\Pi'$ be a deterministic communication protocol and suppose that $\gamma\stackrel{\mathrm{def}}=\Pi'(x_1,x_2)=\Pi'(y_1,y_2)$. Then, $\Pi'(x_1,y_2)=\Pi'(y_1,x_2)=\gamma$.
\end{authorclaim}
In addition, a basic notion of communication complexity that underlies many of its lower bound proofs is that of a fooling set, defdined as follows.
\begin{authorclaim}{Definition2.2 (fooling set (cf., e.g., [7, Sec.1.3] or [10, Chap.1]))}
We say that $S\subseteq\{0,1\}^{n+n}$ is a fooling set for $f:\{0,1\}^{n+n}\to\{0,1\}$ if every $f$-monochromatic rectangle contains at most one point in $S$, where an $f$-monochromatic rectangle is a set $X\times Y$ such that $X,Y\subseteq\{0,1\}^n$ and $f$ is constant on $X\times Y$ (i.e., $f(x,y)=f(x',y')$ for every $(x,y),(x',y')\in X\times Y$).
\end{authorclaim}
Note that a fooling set cannot contain two pre-images of $f^{-1}(0)$ (resp., $f^{-1}(1)$) that differ only on one coordinate; that is, if $(x,y)$ and $(x',y')$ are in a fooling set for $f$ and $f(x,y)=f(x',y')$, then $x\ne x'$ and $y\ne y'$ (because two points that differ on a single coordinate constitute an $f$-monochromatic rectangle).
\subsection*{2.2 Specific Background About Randomness Extraction}
As stated above, an $(\ell,k)$-source is a distribution over $\ell$-bit long strings having min-entropy at least $k$, where the min-entropy of a random variable $X$ is $\min_{v\in\operatorname{Supp}(X)}\{\log_2(1/\Pr[X=v])\}$. That is, $X$ has min-entropy $k$ if and only if for every $v$ it holds that $\Pr[X=v]\le2^{-k}$.
We say that $\mathrm{EXT}:\{0,1\}^d\times\{0,1\}^\ell\to\{0,1\}^m$ is a (seeded) $(k,\epsilon)$-extractor if for every random variable $X$ of min-entropy $k$ the total variation distance between $\mathrm{EXT}(U_d,X)$ and $U_m$ is at most $\epsilon$, where $U_n$ denotes the uniform distribution on $\{0,1\}^n$. In this case $\epsilon$ is called the error of EXT, and $d$ is its seed length.
We say that $\mathrm{EXT}:\{0,1\}^\ell\times\{0,1\}^\ell\to\{0,1\}^m$ is a two-source extractor for independent $(\ell,k)$-sources if for every two independent random variables $X$ and $Y$ of min-entropy $k$ the total variation distance between $\mathrm{EXT}(X,Y)$ and $U_m$ is at most $\epsilon$, called its error. This definition is readily extended to independent sources of parameters $(\ell_1,k_1)$ and $(\ell_2,k_2)$ respectively.
\section*{3 The Public-Randomness Model}
For a protocol $\Pi$ in the public-randomness model, we denote by $\Pi(x,y;r)$ the transcript of the communication on input $(x,y)\in\{0,1\}^{n+n}$ and randomness $r\in\{0,1\}^\ell$. The output of such a protocol is determined by its transcript (e.g., it may be its last bit), and is denoted $\overline\Pi(x,y;r)$.
\begin{authorclaim}{Definition3.1 (communication complexity with a weak public source)}
An $(\ell,k)$-public-randomness protocol for computing a function $f:\{0,1\}^n\times\{0,1\}^n\to\{0,1\}$ is a protocol that satisfies $\Pr[\overline\Pi(x,y;\Xi)=f(x,y)]\ge2/3$, for every $(x,y)\in\{0,1\}^{n+n}$ and every $(\ell,k)$-source $\Xi$.
\end{authorclaim}
\begin{authorclaim}{Theorem3.2 (a general lower bound)}
Suppose that $f:\{0,1\}^{n+n}\to\{0,1\}$ has a fooling get of size $2^m$. Then, any $(\ell,k)$-public-randomness protocol for computing $f$ has communication complexity at least $\min(m-1,\ell-k-1)/2$.
\end{authorclaim}
\begin{proof}
Suppose that $f$ has a $(\ell,k)$-public-randomness protocol, denoted $\Pi$, of communication complexity $t\le(n-1)/2$. We first observe that there exists a dense set of possible source-outcomes $R$ and two input pairs $(x_1,y_1)$ and $(x_2,y_2)$ that reside in the fooling set such that $\Pi$ in constant on all triples $(x_i,y_i,r)$, where $r\in R$ and $i\in\{1,2\}$. The theorem will follow by using the standard ``corners lemma'' (in a non-standard way) and defining a source that is uniform over $R$. Details follow.
The following technical claim has nothing to do with communication complexity; it holds for any function $F:[2^m]\times\{0,1\}^\ell\to\{0,1\}^t$, where in the current case $[2^m]$ represents the indices of the strings in the fooling set (for $f$), $\{0,1\}^\ell$ represents possible outcomes of the public source, and $\{0,1\}^t$ represents possible transcripts of $\Pi$.
\begin{authorclaim}{Claim3.2.1 (a simple combinatorial claim)}
Let $F:[2^m]\times\{0,1\}^\ell\to\{0,1\}^t$. Then, for any $S=\{(x_1,y_1),\ldots,(x_{2^m},y_{2^m})\}$ and $t\le(m-1)/2$, there exist distinct $i,j\in[2^m]$, a string $\gamma\in\{0,1\}^t$, and a set $R\subseteq\{0,1\}^\ell$ of density at least $2^{-2t-1}$ such that for every $r\in R$ it holds that $F(i,r)=F(j,r)=\gamma$.
\end{authorclaim}
We will apply Claim3.2.1 to the hitting set $S$ and to the function $F(i,r)\stackrel{\mathrm{def}}=\Pi(x_i,y_i;r)$. But let us prove the claim first.
\begin{proof}
A simple counting implies that, for every $i\in[2^m]$, there exist $\gamma_i\in\{0,1\}^t$ and a set $R_i\subseteq\{0,1\}^\ell$ of density $2^{-t}$ such that for every $r\in R_i$ it holds that $F(i,r)=\gamma_i$. Similarly, there exist $\gamma\in\{0,1\}^t$ and $G\subseteq[2^m]$ of density $2^{-t}$ such that $\gamma_i=\gamma$ for every $i\in G$.
The key observation is that if $t\le(n-1)/2$, then there exist distinct $i,j\in G$ such that $|R_i\cap R_j|\ge2^{\ell-2t-1}$. This is shown by fixing an arbitrary $G'\subseteq G$ of size $2^{t+1}$, which is possible since $2^{t+1}\le2^{n-t}$, and assuming towards the contradiction that, for every distinct $i,j\in G'$, it holds that $|R_i\cap R_i|<2^{\ell-2t-1}$. Then, we get
\begin{align*}
\left|\bigcup_{i\in G'}R_i\right|&\ge\sum_{i\in G'}|R_i|-\sum_{i\ne j\in G'}|R_i\cap R_j|\\
&>2^{t+1}\cdot2^{\ell-t}-\binom{2^{t+1}}2\cdot2^{\ell-2t-1}\\
&>2^{\ell+1}-2^{2t+1}\cdot2^{\ell-2t-1}=2^\ell,
\end{align*}
which is impossible. The claim follows by fixing $i\ne j$ such that $|R_i\cap R_j|\ge2^{\ell-2t-1}$, and defining $R=R_i\cap R_j$.
\end{proof}
Applying Claim3.2.1 to the hitting set $S=\{(x_1,y_1),\ldots,(x_{2^m},y_{2^m})\}$ of the hypothesis, while letting $F(i,r)\stackrel{\mathrm{def}}=\Pi(x_i,y_i;r)$ and using $t\le(m-1)/2$, we infer that the fooling set contains two points $(x_i,y_i)$ and $(x_j,y_j)$ such that $\Pi(x_i,y_i;r)=\Pi(x_j,y_j;r)=\gamma$ holds for any $r\in R$, where $R\subseteq\{0,1\}^\ell$ has density at least $2^{-2t-1}$.
Next, applying the ``corners lemma'' (i.e., Claim2.1), we infer that $\Pi(x_i,y_i;r)=\Pi(x_i,y_j;r)=\Pi(x_i,y_j;r)=\Pi(x_j,y_j;r)$ for every $r\in R$. Note that this application of the ``corners lemma'' refers to the residual deterministic protocols $\Pi'_r(x,y)=\Pi(x,y;r)$, for all $r\in R$, and it implies that $\Pi'_r(x_i,y_j)=\Pi'_r(x_j,y_i)=\gamma$ for each $r\in R$.
Lastly, picking an $(\ell,\ell-2t-1)$-source that is uniform on $R$, we infer that, when fed with randomness from this source, the execution of $\Pi$ does not distinguish these four input-pairs (i.e., $(x_i,y_i)$, $(x_i,y_j)$, $(x_i,y_j)$ and $(x_j,y_j)$). On the other hand, by hypothesis that $(x_i,y_i)$ and $(x_j,y_j)$ belong to a fooling set, these four input-pairs cannot have the same $f$-value (i.e., it cannot be that $f(x_i,y_i)=f(x_i,y_j)=f(x_j,y_i)=f(x_j,y_j)$, since this would mean that $(x_i,y_i)$ and $(x_j,y_j)$ reside in the $f$-monochromatic rectangle $\{x_i,x_j\}\times\{y_i,y_j\}$). Hence, the hypothesis that $\Pi$ is a $(\ell,k)$-public-randomness protocol for $f$ implies that the foregoing source has min-entropy below $k$; that is, $\ell-2t-1<k$. The theorem follows, since we established $t>(\ell-k-1)/2$, under the hypothesis $t\le(m-1)/2$.
\end{proof}
\subsection*{An archetypical corollary}
Recalling that equality has a fooling set of size $2^n$ and applying Theorem3.2, we get
\begin{authorclaim}{Corollary3.3 (lower bound for equality)}
Any $(\ell,k)$-public-randomness protocol for computing equality of $n$-bit strings has communication complexity at least $\min(n-1,\ell-k-1)/2$.
\end{authorclaim}
This lower bound is tight up to a constant factor, since equality has a constant communication protocol in the standard public-randomness model and the following generic result definitely applies to it.

\begin{authorclaim}{Theorem3.4 (a generic upper bound)}
Suppose that $f:\{0,1\}^{n+n}\to\{0,1\}$ has communication complexity $C^{\mathrm{pub}}(f)$ in the standard public-randomness model. Then, for every $k\le\ell$ such that $k>\log_2n+O(1)$, there exists an $(\ell,k)$-public-randomness protocol for computing $f$ with communication complexity $O(\ell-k)\cdot C^{\mathrm{pub}}(f)$.
\end{authorclaim}
Recall that equality has constant communication complexity in the standard public-randomness model.
\begin{proof}
Recall that the randomness complexity of any protocol for computing $f$ can be reduced to $m\stackrel{\mathrm{def}}=\log_2n+O(1)$ (while possibly increasing its communication complexity by a constant factor).\footnote{See, e.g., [1, Thm.5] and [9]. The basic argument leaves the communication complexity intact, but increases the error probability by an arbitrary small constant, where this constant effects the additive constant in $m$. To regain the original error bound, three repetitions suffice.} The key observation is that the parties can emulate the extraction of $m$ almost-random bits from the public $(\ell,k)$-source, by trying all possible seeds for an adequate randomness extractor, and use the extracted bit to emulate the original randomized protocol. Specifically, for $k\ge m$, such extraction is possible using a (perfectly random) seed of length $d\stackrel{\mathrm{def}}=\log(\ell-k)+O(1)$ (see, e.g., [11, Sec.3.1]). Hence, the parties can emulate the randomized protocol by invoking it $2^d$ times using as randomness the ``extracted outputs'' under all possible seeds. Details follow.
Let $\mathrm{EXT}(s,r)$ denote the output of the extractor $\mathrm{EXT}:\{0,1\}^d\times\{0,1\}^\ell\to\{0,1\}^m$ on seed $s$ and source outcome $r$. Then, given (defective) public-randomness $r\in\{0,1\}^\ell$, the parties emulate $2^d$ invocations of the standard randomized protocol such that in the $i$th invocation they use public-randomness $\mathrm{EXT}(i,r)$, where $i\in[2^d]\equiv\{0,1\}^d$, and rule by majority. Actually, we use a randomized protocol for the standard model that has error probability at most $0.1$ (rather than at most $1/3$), which can be obtained by a constant number of repetitions.
We claim that if EXT has error $0.05$ on any $(\ell,k)$-source $R$, then, for every fixed input pair, with probability at least $2/3$ over the outcome of $R$, the majority of the extracted values (over all possible seeds) yield protocol executions with the correct output. This is the case because otherwise the statistical difference between $\mathrm{EXT}(U_d,R)$ and $U_m$ is at least $\frac13\cdot\frac12-0.1>0.05$, where the first (resp., second) term represents a lower bound (resp., upper bound) on the probability that the protocol yields a wrong answer when run with randomness $\mathrm{EXT}(U_d,R)$ (resp., $U_m$). This yields an $(\ell,k)$-public-randomness protocol of communication complexity $2^d\cdot O(C^{\mathrm{pub}}(f))=O(\ell-k)\cdot C^{\mathrm{pub}}(f)$.
\end{proof}
\subsection*{On the gap between the lower and upper bound}
The bounds provided by Theorems3.3 and3.4 leave a gap of a factor $\Theta(C^{\mathrm{pub}}(f))$ in the non-trivial case (i.e., $\Omega(\ell-k)$ versus $O(\ell-k)\cdot C^{\mathrm{pub}}(f)$). The following example implies that the gap cannot be closed by increasing the lower bound.
\begin{authorclaim}{Proposition3.5 (improved upper bound)}
For every $m<n$, there exists a function $f:\{0,1\}^{n+n}\to\{0,1\}$ that satisfies the following two conditions:
\begin{enumerate}
\item The function $f$ has communication complexity $C^{\mathrm{pub}}(f)=\Theta(m)$ in the standard public-randomness model;
\item For every $k\le\ell$ such that $k>\log_2n+O(1)$, there exists an $(\ell,k)$-public-randomness protocol for computing $f$ with communication complexity $O(\ell-k)+O(C^{\mathrm{pub}}(f))$.
\end{enumerate}
\end{authorclaim}
\begin{proof}
Consider the function $f(x'x'',y'y'')=\mathrm{EQ}(x',y')\oplus\mathrm{IP}_2(x'',y'')$, where $|x''|=m=n-|x'|$, EQ denotes the equality function, and $\mathrm{IP}_2$ denotes inner-product mod2. Then, $C^{\mathrm{pub}}(f)\ge C^{\mathrm{pub}}(\mathrm{IP}_2)=\Omega(m)$, where the first inequality follows by a straightforward reduction and the lower bound is proved in [3]. We obtain an $(\ell,k)$-public-randomness protocol for computing $f$ with communication complexity $O(\ell-k)\cdot C^{\mathrm{pub}}(\mathrm{EQ})+m+1=O(\ell-k)+O(C^{\mathrm{pub}}(f))$, by combining the generic protocol for EQ (see Theorem3.4) with the straightforward deterministic protocol for $\mathrm{IP}_2$.
\end{proof}
\section*{4 The Private-Randomness Model}
For a protocol $\Pi$ in the private-randomness model, we denote by $\Pi((x,r),(y,s))$ the transcript of the communication on input $(x,y)\in\{0,1\}^{n+n}$ with private randomness $r,s\in\{0,1\}^\ell$; that is, the first (resp., second) party gets input $x$ (resp., $y$) and private randomness $r$ (resp., $s$). The output of such a protocol is determined by its transcript (e.g., it may be its last bit), and is denoted $\overline\Pi((x,r),(y,s))$.
\begin{authorclaim}{Definition4.1 (communication complexity with weak private sources)}
An $(\ell,k)$-private-randomness protocol for computing a function $f:\{0,1\}^n\times\{0,1\}^n\to\{0,1\}$ is a protocol that satisfies $\Pr[\overline\Pi((x,\Xi'),(y,\Xi''))=f(x,y)]\ge2/3$, for every $(x,y)\in\{0,1\}^{n+n}$ and every pair of independent $(\ell,k)$-sources $\Xi'$ and $\Xi''$.
\end{authorclaim}
\begin{authorclaim}{Theorem4.2 (a general lower bound)}
Suppose that $f:\{0,1\}^{n+n}\to\{0,1\}$ has a fooling get of size $2^m$. Then, any $(\ell,k)$-private-randomness protocol for computing $f$ has communication complexity at least $\min(m-1,\ell-k-1)/2$.
\end{authorclaim}
\begin{proof}
The proof is analogous to the proof of Theorem3.2. We start with a hypothetical $(\ell,k)$-private-randomness protocol, denoted $\Pi$, that computes $f$ with communication complexity $t\le(n-1)/2$. Then, we apply Claim3.2.1 to the (somewhat less natural) function $F:[2^m]\times\{0,1\}^\ell\to\{0,1\}^t$ defined by $F(i,r)\stackrel{\mathrm{def}}=\Pi((x_i,r),(y_i,r))$, where $S=\{(x_1,y_1),\ldots,(x_{2^m},y_{2^m})\}$ is a fooling set for $f$. Hence, we infer that there exist distinct $i,j\in[2^m]$, a string $\gamma\in\{0,1\}^t$, and a set $R\subseteq\{0,1\}^\ell$ of density at least $2^{-2t-1}$ such that for every $r\in R$ it holds that $\Pi((x_i,r),(y_i,r))=\Pi((x_j,r),(y_j,r))=\gamma$.
Now, applying Claim2.1 thrice, we infer that $\Pi((x_a,r),(y_b,s))=\gamma$ for every $r,s\in R$ and $a,b\in\{i,j\}$. Specifically, for both $a\in\{i,j\}$ and every $r,s\in R$, considering the residual protocol $\Pi'_a(r,s)=\Pi((x_a,r),(y_a,s))$ and using $\Pi((x_a,r),(y_a,r))=\Pi((x_a,s),(y_a,s))=\gamma$, we infer that $\Pi((x_a,r),(y_a,s))=\gamma$. Hence, $\Pi((x_i,r),(y_i,s))=\gamma=\Pi((x_j,r),(y_j,s))$. Now, considering the residual protocol $\Pi'_{r,s}(x,y)=\Pi((x,r),(y,s))$ and using $\Pi((x_i,r),(y_i,s))=\Pi((x_j,r),(y_j,s))$, we get that $\Pi((x_i,r),(y_j,s))=\gamma=\Pi((x_j,r),(y_i,s))$.
Picking a pair of independent $(\ell,\ell-2t-1)$-sources that are each uniform on $R$, we infer that the execution of $\Pi$ does not distinguish the four input-pairs $(x_i,y_i)$, $(x_i,y_j)$, $(x_i,y_j)$ and $(x_j,y_j)$. On the other hand, by hypothesis that $(x_i,y_i)$ and $(x_j,y_j)$ belong to a fooling set, and so these four input-pairs cannot have the same $f$-value. Hence, the hypothesis that $\Pi$ is a $(\ell,k)$-private-randomness protocol for $f$ implies that $\ell-2t-1<k$. The theorem follows, since we established $t>(\ell-k-1)/2$, under the hypothesis $t\le(m-1)/2$.
\end{proof}
\begin{authorclaim}{Theorem4.3 (a generic upper bound)}
Suppose that $f:\{0,1\}^{n+n}\to\{0,1\}$ has communication complexity $C^{\mathrm{priv}}(f)$ in the standard private-randomness model. Then, for every $k\le\ell$ such that $k>2\log_2n+2\log_2\ell+O(1)$, there exists an $(\ell,k)$-private-randomness protocol for computing $f$ with communication complexity $\min(2(\ell-k)+3\log_2n+O(C^{\mathrm{priv}}(f)),\ell+\log_2n+O(C^{\mathrm{priv}}(f))))$.
\end{authorclaim}
\begin{proof}
The bounds follow by having one party send a $\min(2\cdot(\ell-k+\log_2n+O(1)),\ell)$-bit long prefix of its private randomness to the second party, who extracts almost perfect randomness from the two outcomes (using a two-source extractor), sends one half of it back, and then both parties execute the standard protocol. Details follow.
First, recall that the randomness complexity of any protocol for computing $f$ can be reduced to $m\stackrel{\mathrm{def}}=\log_2n+O(1)$ (while possibly increasing its communication complexity by a constant factor). Second, recall that a seedless (two-source) randomness extractor can extract $2m$ almost random bits from an $(\ell,k)$-source and an independent $(\ell',k')$-source, provided that $2m\le k+k'-\max(\ell,\ell')-O(1)$ (see [3, Thm.7(2)]).\footnote{Note that for every $d\ge0$, a $(\ell,k)$-source may be viewed as an $(\ell+d,k)$-source.} Now, if $\ell'\stackrel{\mathrm{def}}=2\cdot(\ell-k+\log_2n)+O(1)\le\ell$, then an $\ell'$-bit prefix of an $(\ell,k)$-source has min-entropy $k'\stackrel{\mathrm{def}}=\ell'-(\ell-k)=(\ell-k)+2\log_2n+O(1)$, and so $k+k'-\max(\ell,\ell')-O(1)=2\log_2n+O(1)$. Hence, sending the prefix of the first source sent to the second party, allows it to extract $2\log_2n+O(1)$ bits that are almost random. Sending half of these bits to the first party allows the two parties to emulate the original protocol. The communication complexity of the proposed protocol is at most $\ell'+\log_2n+O(1)+O(C^{\mathrm{priv}}(f))$, which equals $2(\ell-k)+3\log_2n+O(C^{\mathrm{priv}}(f))$.
As for the case of $\ell'>\ell$, recall that a seedless (two-source) randomness extractor can extract $2m$ almost random bits from a pair of independent $(\ell,k)$-source, provided that $2m\le k-2\log_2\ell-O(1)$ (see [3, Thm.7(1)]). In this case, sending the outcome of the first source to the second party allows for the foregoing emulation, at a total communication cost of $\ell+\log_2n+O(C^{\mathrm{priv}}(f))$.
\end{proof}
